\documentclass[titlepage, a4paper, 12pt]{article}
\usepackage[paper=a4paper, 
            top=2.5cm, 
            bottom=2.5cm, 
            left=2cm, 
            right=2cm, 
            headheight=1.5cm, 
            footskip=1.2cm]{geometry}

\usepackage[utf8]{inputenc}          % Кодировка ввода
\usepackage[T2A]{fontenc}            % Шрифты
\usepackage[english, russian]{babel} % Корректная работа с английским и русским текстом

\usepackage{amsmath, amsfonts, amssymb, amsthm, mathtools} % Файлы американского математического сообщества
\usepackage{hyperref}
\usepackage{enumitem}
\usepackage{tabularx}
\usepackage{subfig}
\usepackage{graphicx}
\graphicspath{{fig/}}

\usepackage{tikz}
\usetikzlibrary{arrows.meta,positioning,shapes.geometric}
\usepackage{adjustbox}

\usepackage{fancyhdr}
\pagestyle{fancy}
\renewcommand{\headrulewidth}{0.4pt}

\lhead{\small Нелинейная динамика и управление нейронными сетями в ИИ}
\chead{}
\rhead{\small \thepage}

\lfoot{}
\cfoot{}
\rfoot{}

\theoremstyle{definition}
\newtheorem*{definition}{Определение}

\theoremstyle{plain}
\newtheorem{theorem}{Теорема}[section]
\newtheorem{lemma}{Лемма}[section]
\newtheorem{corollary}{Следствие}[section]
\newtheorem{proposition}{Предложение}[section]

\theoremstyle{remark}
\newtheorem*{remark}{Замечание}
\newtheorem*{remarks}{Замечания}
\newtheorem*{example}{Пример}

\DeclareMathOperator*{\esssup}{ess\,sup}
\DeclareMathOperator*{\essinf}{ess\,inf}

\usepackage{subfiles}

\begin{document}
\begin{titlepage}
  \title{Нелинейная динамика и управление нейронными сетями в искусственном интеллекте}
  \author{конспект лекций}
  \date{}
  \maketitle
\end{titlepage}

\section*{Предисловие}
Это конспект лекций по курсу <<Нелинейная динамика и управление нейронными сетями в искусственном интеллекте>>, прочитанному в осеннем семестре 2025 года Мокаевым Тимуром Назировичем для бакалавров 4 курса программы <<Прикладная математика, программирование и искусственный интеллект>> и магистров 2 курса <<Математического моделирования, программирования и искусственного интеллекта>> Санкт-Петербургского государственного университета.
Текст был набран студентами по оригинальным презентациям с занятий, изначально подготовленным на английском языке, и исходным статьям (см. список в конце), поэтому может содержать опечатки и не совсем удачные переводы англоязычных терминов, для которых не был найден устоявшийся русский аналог, а также стилистически отличаться от темы к теме\footnote{Для будущих поколений: внимательно читайте как текст, так и сходный код, который вам сгенерировала LLM!}.
\TeX-файлы находятся в открытом доступе на \hyperref{https://github.com/IKurkov/Nonlinear_dynamics_and_control_of_neural_networks_in_AI}{}{}{GitHub}{}.

\tableofcontents
\newpage
\section{Биологические и искусственные нейроны. Исторический контекст искусственных нейронных сетей}
\subfile{lections/lec1_1}
\subfile{lections/lec1_2}
\newpage

\section{Функциональные пространства и универсальная аппроксимация}
\subfile{lections/lec2_part1_1}
\subfile{lections/lec2_part1_2}
\subfile{lections/lec2_part1_3}
\newpage

\section{Классические теоремы об универсальной аппроксимации для функций в \textit{C} и $L^p$}
\subfile{lections/lec2_part2_1}
\subfile{lections/lec2_part2_2}
\newpage

\section{Универсальная аппроксимация нелинейных операторов}
\subfile{lections/lec3_1}
\subfile{lections/lec3_2_UAT_for_functionals}
\subfile{lections/lec3_3}
\newpage

\section{Generalized Translation Networks и их применение к задачам идентификации и управления}
\subfile{lections/lec4_1_GTN}
\subfile{lections/lec4_2_4models}
\newpage

\section{Функциональный многослойный перцептрон (FMLP)}
\subfile{lections/lec5_1}
\subfile{lections/lec5_2}
\newpage

\section{Deep Operator Networks (DeepONets) и роль глубины и ширины}
\subfile{lections/lec6_1}
\subfile{lections/lec6_2}
\newpage

\section{Глубокое обучение через призму динамических систем}
\subfile{lections/lec7}
\newpage

\section{Универсальная аппроксимация свёрточными сетями}
\subfile{lections/lec8_1}
\subfile{lections/lec8_2}
\newpage

\section{Свойство универсальной аппроксимации}
\subfile{lections/lec9_1}
\subfile{lections/lec9_2}
\newpage

\addcontentsline{toc}{section}{Список литературы}
\bibliographystyle{unsrt}
\bibliography{references}

\end{document}